% !TEX TS-program = pdflatex
\documentclass[11pt,a4paper]{report}

% ---------------------------------------------------------------------------
% Encodage, langue, police
% ---------------------------------------------------------------------------
\usepackage[utf8]{inputenc}
\usepackage[T1]{fontenc}
\usepackage[french]{babel}
\usepackage{lmodern}
\usepackage{microtype}

% ---------------------------------------------------------------------------
% Mise en page
% ---------------------------------------------------------------------------
\usepackage[margin=2.4cm,top=2.8cm,bottom=2.6cm]{geometry}
\usepackage{fancyhdr}
\usepackage{titlesec}
\usepackage{setspace}
\onehalfspacing

\pagestyle{fancy}
\fancyhf{}
\fancyhead[L]{\small\itshape Cloud Computing \& Déploiement IA --- TP 2}
\fancyhead[R]{\small\itshape Conteneuriser un modèle et optimiser son image}
\fancyfoot[C]{\small\thepage}
\renewcommand{\headrulewidth}{0.4pt}

\titleformat{\chapter}[display]{\normalfont\huge\bfseries}{\chaptertitlename~\thechapter}{0.6ex}{\titlerule}
\titlespacing*{\chapter}{0pt}{-20pt}{18pt}

% ---------------------------------------------------------------------------
% Tableaux, figures, listes
% ---------------------------------------------------------------------------
\usepackage{booktabs}
\usepackage{longtable}
\usepackage{array}
\usepackage{xcolor}
\usepackage{colortbl}
\usepackage{graphicx}
\usepackage{float}
\usepackage{caption}
\usepackage{remreset}
\usepackage{enumitem}
\usepackage{amsmath}
\usepackage{xurl}
\usepackage[most]{tcolorbox}
\usepackage{listings}
\usepackage{hyperref}

\definecolor{accent}{HTML}{1F4E79}
\definecolor{alert}{HTML}{A61C1C}
\definecolor{okgreen}{HTML}{1E6B34}
\definecolor{codebg}{HTML}{F5F6F8}
\definecolor{notebg}{HTML}{EEF3F9}
\definecolor{warnbg}{HTML}{FBEDED}

\hypersetup{
  colorlinks=true,
  linkcolor=accent,
  urlcolor=accent,
  citecolor=accent,
  pdftitle={Rapport TP 2 - Conteneuriser un modele et optimiser son image},
  pdfauthor={TP2 - Cloud Computing & Deploiement IA}
}

% babel-francais met deja les libelles de legende en petites capitales ;
% on ne force pas le gras, Latin Modern n'a pas de petites capitales grasses.
\captionsetup{font=small,labelfont={color=accent},labelsep=endash}

% Numerotation continue des figures et tableaux (1, 2, 3...) et non
% <chapitre>.<n>, qui se confondrait avec la numerotation des sections.
% report remet figure et table a zero a chaque \chapter ; on retire ce reset.
% On utilise remreset et non chngcntr, car le paquet caption recharge le
% reset apres le passage de chngcntr et la numerotation repart de 1.
\renewcommand{\thefigure}{\arabic{figure}}
\renewcommand{\thetable}{\arabic{table}}
\makeatletter
\@removefromreset{figure}{chapter}
\@removefromreset{table}{chapter}
\makeatother
\setlist{itemsep=2pt,topsep=4pt}
\setlist[enumerate]{label=\arabic*.,leftmargin=*}

% ---------------------------------------------------------------------------
% Listings
% ---------------------------------------------------------------------------
\lstdefinelanguage{Dockerfile}{
  keywords={FROM,RUN,CMD,LABEL,MAINTAINER,EXPOSE,ENV,ADD,COPY,ENTRYPOINT,VOLUME,USER,WORKDIR,ARG,ONBUILD,STOPSIGNAL,HEALTHCHECK,SHELL,AS},
  sensitive=false,
  comment=[l]{\#},
  morestring=[b]"
}
\lstdefinelanguage{Yaml}{
  keywords={services,volumes,build,image,container_name,ports,depends_on,condition,service_completed_successfully,restart,environment,healthcheck,test,interval,timeout,start_period,retries,command,entrypoint},
  sensitive=false,
  comment=[l]\#,
  morestring=[b]"
}
\lstset{
  basicstyle=\ttfamily\scriptsize,
  backgroundcolor=\color{codebg},
  frame=single,
  rulecolor=\color{black!25},
  framesep=4pt,
  breaklines=true,
  breakatwhitespace=false,
  columns=fullflexible,
  keepspaces=true,
  showstringspaces=false,
  numbers=left,
  numbersep=6pt,
  numberstyle=\tiny\color{black!45},
  keywordstyle=\color{accent}\bfseries,
  commentstyle=\color{black!55}\itshape,
  stringstyle=\color{okgreen},
  tabsize=4,
  upquote=true,
  extendedchars=true,
  escapeinside={(*@}{@*)},
  literate={é}{{\'e}}1 {è}{{\`e}}1 {ê}{{\^e}}1 {à}{{\`a}}1
           {û}{{\^u}}1 {ç}{{\c c}}1 {î}{{\^i}}1 {ô}{{\^o}}1
           {ù}{{\`u}}1 {œ}{{\oe}}1 {É}{{\'E}}1,
}
\lstset{language=Dockerfile}

% Encadres
\newtcolorbox{encadre}[1][]{%
  enhanced,breakable,colback=notebg,colframe=accent,boxrule=0.6pt,
  left=6pt,right=6pt,top=5pt,bottom=5pt,arc=1.5pt,
  fonttitle=\bfseries\small,title={#1}}
\newtcolorbox{alerte}[1][]{%
  enhanced,breakable,colback=warnbg,colframe=alert,boxrule=0.6pt,
  left=6pt,right=6pt,top=5pt,bottom=5pt,arc=1.5pt,
  fonttitle=\bfseries\small,title={#1}}
\newtcolorbox{resultat}[1][]{%
  enhanced,breakable,colback=white,colframe=okgreen,boxrule=0.8pt,
  left=6pt,right=6pt,top=5pt,bottom=5pt,arc=1.5pt,
  fonttitle=\bfseries\small,title={#1}}

% Macro : cellule centree
\newcommand{\cc}[1]{\centering #1}
\newcommand{\num}[1]{\textbf{#1}}

% Macro : capture d'ecran du dossier screens/, inseree dans le corps du
% rapport a l'emplacement exact de l'etape qu'elle documente.
%   #1 = nom de fichier (les espaces sont acceptes par graphicx)
%   #2 = label de reference
%   #3 = legende
\newcommand{\capture}[3]{%
  \begin{figure}[H]
  \centering
  \includegraphics[width=0.95\textwidth,height=0.52\textheight,keepaspectratio]{#1}
  \caption{#3}
  \label{#2}
  \end{figure}}

% Variante etroite, pour les captures tres panoramiques (ratio > 3).
\newcommand{\capturepan}[3]{%
  \begin{figure}[H]
  \centering
  \includegraphics[width=0.74\textwidth,height=0.34\textheight,keepaspectratio]{#1}
  \caption{#3}
  \label{#2}
  \end{figure}}

% ---------------------------------------------------------------------------
\begin{document}

% ============================ PAGE DE GARDE ===============================
\begin{titlepage}
\thispagestyle{empty}
\centering

\vspace*{2.5cm}

{\large université mundiapolis\par}
\vspace{0.3cm}
{\large Cloud Computing \& Déploiement IA\par}

\vspace{2.5cm}

{\color{accent}\rule{\textwidth}{0.8pt}\par}

\vspace{0.8cm}

{\Huge\bfseries TP 2\par}

\vspace{0.5cm}

{\LARGE\bfseries Conteneuriser un modèle et optimiser son image\par}

\vspace{0.8cm}

{\color{accent}\rule{\textwidth}{0.8pt}\par}

\vfill

\begin{minipage}{0.7\textwidth}
\centering
\begin{tabular}{@{}ll@{}}
\textbf{Présenté par} & Aymane El yamani \\
\textbf{Encadré par} & Pr. Mohammed AMEKSA \\
\textbf{Année universitaire} & 2026--2027 \\
\end{tabular}
\end{minipage}

\vfill


\vspace{1.5cm}

\end{titlepage}

% ============================ SOMMAIRE ====================================
\tableofcontents
\clearpage
\listoftables
\clearpage
\listoffigures
\newpage

% ============================ CHAPITRE 1 =================================
\chapter{Introduction}

\section{Contexte}

Ce rapport présente le travail réalisé dans le cadre du TP 2 du cours
\textbf{Cloud Computing \& Déploiement IA}.

Dans ce TP, nous travaillons sur un modèle de classification des revenus basé
sur le jeu de données \textbf{Adult Income (UCI)} et développé avec
\textbf{scikit-learn}. L'objectif est de conteneuriser le modèle avec
\textbf{Docker}, puis d'optimiser et de sécuriser son image.

Les objectifs principaux du TP sont :

\begin{itemize}
    \item conteneuriser l'entraînement et l'inférence du modèle ;
    \item optimiser la taille de l'image Docker et mesurer les temps de construction ;
    \item utiliser \textbf{Docker Compose} pour gérer plusieurs conteneurs avec un volume partagé ;
    \item analyser les vulnérabilités de l'image avec \textbf{Trivy} et les corriger.
\end{itemize}

\section{Prérequis et environnement de travail}

\begin{table}[H]
\centering\small
\begin{tabular}{@{}p{4.6cm}p{9.4cm}@{}}
\toprule
\textbf{Élément} & \textbf{Vérification / valeur} \\
\midrule
Docker & 28.5.1 (build e180ab8) \\
Docker Compose & v2 (plugin) \\
Trivy & installé, base de règles \texttt{debian 13.7} \\
Base naive & \texttt{python:3.11} \\
Base optimisée & \texttt{python:3.11-slim} \\
Python local (hors conteneur) & 3.12.10 \\
OS hôte & Windows 11, Docker Desktop (moteur Linux) \\
Données & \texttt{data/adult.csv} --- 32\,561 lignes (UCI Adult réel) \\
\end{tabular}
\caption{Environnement de travail et prérequis vérifiés.}
\end{table}

% ============================ CHAPITRE 2 =================================
\chapter{Le projet : chaîne de traitement et implémentation}

\section{Arborescence du dépôt}

Le dépôt respecte l'arborescence imposée par le support :

\begin{table}[H]
\centering\small
\begin{tabular}{@{}p{5.4cm}p{8.6cm}@{}}
\toprule
\textbf{Fichier / répertoire} & \textbf{Rôle} \\
\midrule
\texttt{src/train.py}     & Entraînement, évaluation, sérialisation du pipeline \\
\texttt{src/predict.py}   & Inférence : mode lot (CLI) et mode service (FastAPI) \\
\texttt{src/common.py}    & \texttt{ColumnTransformer} partagé entre les deux \\
\texttt{data/download.py} & Récupération UCI ou génération hors ligne \\
\texttt{requirements.txt} & Dépendances Python épinglées \\
\texttt{Dockerfile.train} & Image d'entraînement \\
\texttt{Dockerfile.serve} & Image d'inférence naïve \\
\texttt{Dockerfile.serve.opt1 \dots opt6} & Une optimisation par fichier \\
\texttt{Dockerfile.serve.opt} & Image finale \\
\texttt{benchmark.sh}     & Mesure automatisée des trois indicateurs \\
\texttt{.dockerignore}    & Réduction du contexte de construction \\
\texttt{mesure-results.txt} & Relevé des 8 mesures \\
\texttt{trivy.checks.txt} & Rapport Trivy de l'image finale \\
\texttt{screens/}         & 13 captures de démonstration \\
\bottomrule
\end{tabular}
\caption{Arborescence du dépôt et rôle de chaque fichier.}
\end{table}

Le point d'architecture le plus important est le découpage de
\texttt{src/common.py}. Comme le souligne le support, ce module est importé
\emph{à la fois} par l'entraînement et par l'inférence :
« c'est ce qui évite l'écart entre la préparation des données à l'entraînement
et en production ». Concrètement, le prétraitement (imputation, mise à
l'échelle, encodage one-hot) n'est pas écrit deux fois : il est \emph{objet de
pipeline}, donc sérialisé \emph{avec} le modèle dans le fichier
\texttt{model.joblib}. Une transformation de feature introduite en
production sans être répliquée à l'entraînement est une source classique de
\emph{drift} silencieux ; le \texttt{ColumnTransformer} partagé l'élimine par
construction.

\section{Le jeu de données}

Le script \texttt{data/download.py} tente un téléchargement depuis l'archive
UCI et bascule automatiquement sur un générateur synthétique en cas d'échec
réseau, ce qui rend le TP réalisable en salle isolée. Sur cette machine, le
téléchargement a abouti : \texttt{data/adult.csv} contient \textbf{32\,561
lignes} correspondant au jeu UCI \emph{Adult Income} authentique, avec la
distribution de référence :

\begin{table}[H]
\centering
\begin{tabular}{@{}lrr@{}}
\toprule
\textbf{Classe} & \textbf{Effectif} & \textbf{Fréquence} \\
\midrule
\texttt{<=50K} & 24\,720 & 75,92\,\% \\
\texttt{>50K}  & 7\,841  & 24,08\,\% \\
\midrule
Total        & 32\,561 & 100\,\% \\
\bottomrule
\end{tabular}
\caption{Distribution des classes du jeu Adult Income.}
\end{table}

Le déséquilibre de classe (24\,\% de positifs) est un point important pour la
choix du modèle et des métriques : l'exactitude seule est un indicateur
trompeur, ce qui justifie le recours simultané au F1 et à l'AUC.

\section{Entraînement : \texttt{src/train.py}}

Le découpage est \emph{stratifié} (\texttt{stratify=y}) et la graine est fixée
à 42, ce qui rend l'expérience reproductible. Avec \texttt{test\_size=0.2},
on obtient 26\,048 Training / 6\,513 test, ce que confirme
\texttt{metrics.json} (\texttt{n\_entrainement = 26048}).

Le pipeline construit par \texttt{construire\_pipeline()} est le suivant :

\begin{table}[H]
\centering\small
\begin{tabular}{@{}p{3.2cm}p{12.2cm}@{}}
\toprule
\textbf{Bloc} & \textbf{Contenu} \\
\midrule
Prétraitement num. & \texttt{SimpleImputer(median)} puis \texttt{StandardScaler} sur 5 colonnes \\
Prétraitement cat.  & \texttt{SimpleImputer(most\_frequent)} puis
\newline \texttt{OneHotEncoder(handle\_unknown="ignore",}\newline
\texttt{sparse\_output=False)} sur 5 colonnes \\
Modèle            & \texttt{HistGradientBoostingClassifier}\newline
\texttt{(max\_iter=200, random\_state=42)} \\
\bottomrule
\end{tabular}
\caption{Contenu du pipeline de prétraitement et de modélisation.}
\end{table}

Les métriques calculées sont l'exactitude, le F1 et l'AUC, avec un seuil de
décision fixe à 0,5. Résultats obtenus (\texttt{artifacts/metrics.json}) :

\begin{resultat}[Performance du modèle]
\begin{center}
\begin{tabular}{@{}lr@{}}
\toprule
\textbf{Métrique} & \textbf{Valeur} \\
\midrule
Exactitude (accuracy)     & \num{0,8773} \\
F1 (classe positive)      & \num{0,7263} \\
AUC-ROC                   & \num{0,9298} \\
Durée d'entraînement      & 8,24 s \\
Individus en entraînement & 26\,048 \\
Graine                    & 42 \\
\bottomrule
\end{tabular}
\end{center}
\end{resultat}

L'AUC de 0,93 indique un excellent pouvoir de discrimination. L'exactitude de
0,877 est en revanche \emph{modeste} pour Adult Income : les références
publiées dépassent 0,90. L'explication est identifiable dans
\texttt{src/common.py} : seules 10 variables sur 15 sont retenues
(\texttt{fnlwgt}, \texttt{education}, \texttt{race} et
\texttt{native\_country} sont écartées). \texttt{education\_num} est un
proxy partiel de \texttt{education}, mais \texttt{fnlwgt} est la variable la
plus informative du jeu et son absence pénalise sensiblement l'exactitude.
Le F1 de 0,726 est cohérent avec la difficulté de la classe positive
(\texttt{>50K}), minoritaire à 24\,\%.

\section{Inférence : \texttt{src/predict.py}}

Le module offre deux modes, comme le demande le support :

\begin{itemize}
  \item \textbf{mode lot} : \texttt{python -m src.predict --input
    data/sample.csv}, qui affiche les probabilités des 10 premières lignes ;
  \item \textbf{mode service} : \texttt{uvicorn src.predict:app}, qui expose
        l'API FastAPI.
\end{itemize}

L'API comporte trois points d'entrée, et cette tripartition est
particulièrement bien conçue pour le TP 3 :

\begin{table}[H]
\centering\small
\begin{tabular}{@{}p{2.6cm}p{9.4cm}p{2.6cm}@{}}
\toprule
\textbf{Route} & \textbf{Sémantique} & \textbf{Équivalent K8s} \\
\midrule
\texttt{GET /health}     & Sonde de vivacité : le processus répond-il~? & \texttt{livenessProbe} \\
\texttt{GET /ready}      & Sonde de disponibilité : le modèle est-il chargé~? & \texttt{readinessProbe} \\
\texttt{POST /v1/predict} & Prédiction sur un individu validé par Pydantic & Service applicatif \\
\bottomrule
\end{tabular}
\caption{Points d'entrée de l'API et correspondance Kubernetes.}
\end{table}

Le chargement du modèle est \emph{paresseux} (lazy loading) : le
\texttt{joblib} n'est lu qu'au premier appel. Le commentaire du code en
reconnaît lui-même la limite pour la production :
« En production on préfère le charger au démarrage (événement \texttt{lifespan})
pour que la première requête ne paie pas le coût. »

Deux points de qualité sont à souligner positivement. D'abord, le modèle
Pydantic \texttt{Individu} valide les plages numériques
(\texttt{age} entre 17 et 100, \texttt{education\_num} entre 1 et 16,
\texttt{hours\_per\_week} entre 1 et 99) : l'API rejette donc les requêtes
malformées au lieu de produire une prédiction aberrante. Ensuite, la réponse
inclut la \texttt{latence\_ms} mesurée autour de l'appel
\texttt{predict\_proba}, ce qui fournit une métrique exploitable en
supervision.

\section{Le script de mesure : \texttt{benchmark.sh}}
\label{sec:benchmark}

Le script implémente exactement le protocole décrit dans le support et
applique une sequencing méthodologiquement importante : il mesure
\emph{d'abord} la construction à froid avec \texttt{--no-cache}, puis
\emph{ensuite} la reconstruction après modification. L'ordre inverse aurait été
faux : sans le passage préalable en \texttt{--no-cache}, le cache serait déjà
partiellement peuplé et la mesure « à froid » n'aurait plus de sens.

\begin{lstlisting}[language=bash,numbers=left]
T_FROID=$(( $(date +%s) - T0 ))          # docker build --no-cache
echo "# touche par benchmark.sh $(date +%s)" >> src/predict.py
T0=$(date +%s)
docker build -f "${DOCKERFILE}" -t "${IMAGE}" . >/dev/null
T_CHAUD=$(( $(date +%s) - T0 ))          # reconstruction
sed -i '$ d' src/predict.py              # nettoyage
TAILLE=$(docker image inspect "${IMAGE}" --format '{{.Size}}')
\end{lstlisting}

\section{Relevé des mesures}
\label{sec:releve}

Le fichier \texttt{mesure-results.txt} contient les huit mesures. Elles sont
reportées et commentées dans les sections suivantes. La dernière entrée,
\texttt{adult-serve:opt}, correspond à l'image finale et sert de référence
absolute.

\begin{encadre}[Comment lire les sections qui suivent]
Chaque étape de ce chapitre suit le même plan~: la commande, le
\texttt{Dockerfile} correspondant, la mesure relevée, puis la \textbf{capture
d'écran de la manipulation}. Les figures sont donc insérées à l'emplacement
exact de l'étape qu'elles documentent~: il n'y a pas d'annexe de captures à
consulter séparément, et la numérotation des figures suit l'ordre des
étapes du TP.
\end{encadre}

% ============================ CHAPITRE 3 =================================
\chapter{Axe I --- Conteneuriser un modèle IA}

\section{Version naïve : le point de référence}
\label{sec:naif}

\subsection{Construction de l'image d'entraînement}

L'image naïve est construite sur \texttt{python:3.11}, l'image officielle
complète, qui embarque l'ensemble de la chaîne d'outils Debian~: compilateur,
documentation, outils système. \texttt{Dockerfile.train} installe les
dépendances puis lance \texttt{python -m src.train}~:

\begin{lstlisting}[language=Dockerfile,numbers=left]
FROM python:3.11
WORKDIR /app
COPY requirements.txt .
RUN pip install -r requirements.txt
COPY src/ ./src/
COPY data/adult.csv ./data/adult.csv
RUN mkdir -p /artifacts
CMD ["python", "-m", "src.train",
     "--data", "data/adult.csv",
     "--out", "/artifacts/model.joblib"]
\end{lstlisting}

\capture{screens/01.build_adult-train-naif.JPG}{fig:build-train}{Image d'entraînement naïve}

\subsection{Construction de l'image d'inférence}

Le \texttt{Dockerfile.serve} charge \texttt{/artifacts/model.joblib} et expose
la prédiction en ligne de commande. C'est également la version de référence
\emph{sans aucune}~ optimisation~:

\begin{lstlisting}[language=Dockerfile,numbers=left]
FROM python:3.11
WORKDIR /app
COPY requirements.txt .
COPY src/ ./src/
COPY data/sample.csv ./data/sample.csv
RUN pip install -r requirements.txt
ENV MODEL_PATH=/artifacts/model.joblib
ENTRYPOINT ["python", "-m", "src.predict"]
CMD ["--input", "data/sample.csv"]
\end{lstlisting}

\capture{screens/02.build_adult-serve-naif.JPG}{fig:build-serve}{Image d'inférence naïve}

\begin{alerte}[Défaut structurel de la version naïve]
L'ordre des instructions est le point critique. \texttt{COPY src/} est exécuté
\emph{avant} \texttt{RUN pip install}. Or le cache de construction de Docker
invalide une couche dès que \emph{l'instruction ou ses entrées} changent. Toute
modification d'une ligne de \texttt{src/predict.py} invalide donc la couche
\texttt{COPY src/} \emph{et toutes les couches suivantes}, ce qui impose de
rejouer l'installation complète de \texttt{numpy}, \texttt{scipy},
\texttt{pandas} et \texttt{scikit-learn} --- des centaines de mégaoctets à
télécharger et extraire. C'est la réponse à la question 1 du support, et
elle est mesurable~: \num{106 s} de reconstruction pour le naïf.
\end{alerte}

Remarquons que \texttt{Dockerfile.train} ne souffre pas de ce défaut~: il
respecte déjà l'ordre \texttt{requirements.txt} $\rightarrow$
\texttt{pip install} $\rightarrow$ \texttt{src/}. L'antipattern n'a donc été
introduit que dans l'image de service.

\subsection{Sauvegarde des artefacts et première mesure}

Le modèle, les métriques et l'échantillon de test sont conservés hors du cycle
de vie du conteneur par un montage de volume~: \texttt{model.joblib} (450~Ko) et
\texttt{metrics.json} sortent du conteneur, qui peut être détruit sans perdre le
modèle. L'image d'inférence lira le modèle par ce même volume (\S~\ref{sec:compose}).

\begin{lstlisting}[language=bash,numbers=left]
docker build -f Dockerfile.train -t adult-train:naif .
docker run --rm -v $(pwd)/artifacts:/artifacts adult-train:naif
ls -lh artifacts/model.joblib
bash benchmark.sh naif Dockerfile.serve      # -> 624 Mo
\end{lstlisting}

\capture{screens/03.sauvgard_dans_actefact.JPG}{fig:artefacts}{Sauvegarde des artefacts}

\capture{screens/04.mesur la taille naive.JPG}{fig:mesure-naif}{Mesure de l'image naïve}

\section{Optimisation 1 : base \texttt{python:3.11-slim}}
\label{sec:opt1}

\begin{lstlisting}[language=Dockerfile,numbers=left,numbers=none]
FROM python:3.11-slim          # au lieu de python:3.11
\end{lstlisting}

Le support justifie le choix de \texttt{slim} plutôt qu'\texttt{alpine}~: avec
la libc musl, \texttt{numpy} et \texttt{scikit-learn} doivent être recompilés
depuis les sources, ce qui allongerait massivement la construction. Ce conseil
est correct et se vérifie dans la littérature~: les wheels \emph{manylinux}
sont construits contre glibc, pas musl.

\capture{screens/05.optimisation 1.JPG}{fig:opt1}{Optimisation 1 --- base \texttt{slim}}

\section{Optimisation 2 : ordonnancement des \texttt{COPY}}
\label{sec:opt2}

\begin{lstlisting}[language=Dockerfile,numbers=left]
FROM python:3.11-slim
WORKDIR /app
COPY requirements.txt .              # <- isole la couche "dépendances"
RUN pip install -r requirements.txt   # <- le cache survit aux modifications de src/
COPY src/ ./src/                     # <- invalidé souvent, mais peu coûteux
COPY data/sample.csv ./data/sample.csv
ENV MODEL_PATH=/artifacts/model.joblib
ENTRYPOINT ["python", "-m", "src.predict"]
CMD ["--input", "data/sample.csv"]
\end{lstlisting}

C'est l'optimisation la plus importante du TP en termes de gain sur le temps
d'itération. Elle ne change \emph{rien} à la taille de l'image mais divise le
temps de reconstruction par près de neuf.

\capture{screens/06.optimisation 2.JPG}{fig:opt2}{Optimisation 2 --- ordre des \texttt{COPY}}

\section{Optimisation 3 : \texttt{.dockerignore}}
\label{sec:opt3}

Le fichier \texttt{.dockerignore} exclut du contexte de construction les
répertoires et fichiers non nécessaires : \texttt{venv/}, \texttt{.git/},
\texttt{\_\_pycache\_\_/}, \texttt{artifacts/}, \texttt{*.pdf}, etc.

\begin{alerte}[Précision indispensable sur cette optimisation]
Un \texttt{.dockerignore} \textbf{ne modifie jamais la taille de l'image}. Il
réduit le \emph{contexte de construction}, c'est-à-dire la quantité de données
transmise au démon Docker. Son effet est donc visible sur le temps de
construction (transfert du contexte) et sur le risque de fuite de secrets, mais
pas sur \texttt{docker images}. Le support, qui présente cette optimisation
dans un tableau mêlant taille et durées, peut donner l'impression inverse ; les
mesures le confirment : la taille reste rigoureusement identique à celle de
l'optimisation 2 (\num{263 Mo} avant comme après).
\end{alerte}

\capture{screens/07.optimisation 3.JPG}{fig:opt3}{Optimisation 3 --- \texttt{.dockerignore}}

\section{Optimisation 4 : \texttt{pip install -{}-no-cache-dir}}
\label{sec:opt4}

\begin{lstlisting}[language=Dockerfile,numbers=left,numbers=none]
RUN pip install --no-cache-dir -r requirements.txt
\end{lstlisting}

Par défaut, \texttt{pip} conserve les wheels téléchargés dans
\texttt{\~/.cache/pip}. Comme ces fichiers sont créés \emph{dans} la couche
\texttt{RUN}, ils sont figés dans l'image et jamais nettoyés : c'est du poids
mort. L'option \texttt{-{}-no-cache-dir} supprime cette couche.

C'est la deuxième optimisation la plus rentable en taille, et son effet est
aisément isolable puisque les optimisations 1 à 3 n'ont rien modifié à la
taille : \textbf{la taille passe de \num{263 Mo} à \num{172 Mo}, soit
-\num{91 Mo}}. Ce chiffre correspond exactement au poids du cache pip
figé dans l'image.

\capture{screens/08.optimisation 4.JPG}{fig:opt4}{Optimisation 4 --- \texttt{-{}-no-cache-dir}}

\section{Optimisation 5 : construction multi-étapes}
\label{sec:opt5}

\begin{lstlisting}[language=Dockerfile,numbers=left]
FROM python:3.11-slim AS builder
WORKDIR /build
COPY requirements.txt .
RUN python -m venv /opt/venv
RUN /opt/venv/bin/pip install --no-cache-dir -r requirements.txt

FROM python:3.11-slim
WORKDIR /app
ENV PATH="/opt/venv/bin:$PATH"
COPY --from=builder /opt/venv /opt/venv   # seul le résultat est livré
COPY src/ ./src/
COPY data/sample.csv ./data/sample.csv
ENTRYPOINT ["python", "-m", "src.predict"]
CMD ["--input", "data/sample.csv"]
\end{lstlisting}

L'image \emph{builder} contient les outils de compilation (gcc, headers) qui
servent éventuellement à construire des dépendances depuis les sources. L'image
finale n'en copie que le résultat via \texttt{COPY --from=builder}, de sorte
que la chaîne d'outils n'est jamais livrée.

\capture{screens/09.optimisation 5.JPG}{fig:opt5}{Optimisation 5 --- multi-étapes}

\section{Optimisation 6 : utilisateur non-root et \texttt{HEALTHCHECK}}
\label{sec:opt6}

\begin{lstlisting}[language=Dockerfile,numbers=left]
RUN useradd --create-home --uid 10001 appuser \
    && mkdir -p /artifacts \
    && chown appuser:appuser /artifacts
COPY --chown=appuser:appuser src/ ./src/
COPY --chown=appuser:appuser data/sample.csv ./data/sample.csv
USER appuser
EXPOSE 8000
HEALTHCHECK --interval=15s --timeout=5s --start-period=20s --retries=3 \
  CMD ["python", "-c", "import urllib.request; \
        urllib.request.urlopen('http://127.0.0.1:8000/health', timeout=3)"]
ENTRYPOINT ["uvicorn", "src.predict:app", "--host", "0.0.0.0", "--port", "8000"]
\end{lstlisting}

C'est l'optimisation la plus structurante pour le TP 3. Elle embarque quatre
bonnes pratiques :

\begin{enumerate}
  \item \textbf{Utilisateur non-root} : le conteneur s'exécute sous
        \texttt{appuser} (UID 10001) et non sous \texttt{root}. Cela réduit
        fortement la surface d'attaque et correspond au
        \texttt{securityContext} standard de Kubernetes ;
  \item \textbf{UID fixe} (10001) plutôt qu'un UID alloué dynamiquement : cela
        évite toute collision avec un utilisateur existant sur le nœud hôte ;
  \item \textbf{Ownership explicite par \texttt{--chown}} : les fichiers sont
        lisibles par \texttt{appuser}, sans recours à un \texttt{chmod}
        global, qui créerait une couche supplémentaire ;
  \item \textbf{\texttt{ENTRYPOINT} en forme \emph{exec}} (forme tableau JSON)
        et non en forme shell : le processus \texttt{uvicorn} devient le
        PID~1, ce qui lui permet de recevoir directement les signaux
        \texttt{SIGTERM} de l'orchestrateur --- condition indispensable pour
        un arrêt propre sur Kubernetes.
\end{enumerate}

\capture{screens/10.optimisation 6.JPG}{fig:opt6}{Optimisation 6 --- non-root et \texttt{HEALTHCHECK}}

\section{Tableau de synthèse des mesures}
\label{sec:synthese}

Le tableau ci-dessous consolide les huit mesures de
\texttt{mesure-results.txt}. La colonne \emph{delta taille} donne l'écart par
rapport à l'optimisation précédente ; la colonne \emph{delta recons.} idem pour
le temps de reconstruction. La colonne de droite renvoie à la capture
documentant chaque étape.

\begin{table}[H]
\centering\small
\begin{tabular}{@{}r p{3.9cm} rrr rr c@{}}
\toprule
\textbf{\#} & \textbf{Optimisation} & \textbf{Taille} & \textbf{Build à froid} & \textbf{Rebuild} & \textbf{$\Delta$ taille} & \textbf{$\Delta$ recons.} & \textbf{Capture} \\
\midrule
0 & Version naïve (référence)      & \num{624 Mo} & \num{112 s} & \num{106 s} & ---      & --- & \hyperref[fig:mesure-naif]{\ref{fig:mesure-naif}} \\
1 & Base \texttt{python:3.11-slim} & \num{263 Mo} & \num{108 s} & \num{130 s} & $-361$  & $+24$ & \hyperref[fig:opt1]{\ref{fig:opt1}} \\
2 & \texttt{COPY req.} avant \texttt{src/} & \num{263 Mo} & \num{159 s} & \num{15 s} & $0$   & $-115$ & \hyperref[fig:opt2]{\ref{fig:opt2}} \\
3 & \texttt{.dockerignore} complet & \num{263 Mo} & \num{678 s} & \num{5 s}  & $0$     & $-10$ & \hyperref[fig:opt3]{\ref{fig:opt3}} \\
4 & \texttt{pip install -{}-no-cache-dir} & \num{172 Mo} & \num{154 s} & \num{16 s} & $-91$ & $+11$ & \hyperref[fig:opt4]{\ref{fig:opt4}} \\
5 & Construction multi-étapes     & \num{176 Mo} & \num{208 s} & \num{18 s} & $+4$   & $+2$ & \hyperref[fig:opt5]{\ref{fig:opt5}} \\
6 & Non-root + \texttt{HEALTHCHECK} & \num{176 Mo} & \num{253 s} & \num{30 s} & $0$    & $+12$ & \hyperref[fig:opt6]{\ref{fig:opt6}} \\
\midrule
& \textbf{Image finale \texttt{adult-serve:opt}} & \num{176 Mo} & \num{171 s} & \num{14 s} & & & \hyperref[fig:dockerimages]{\ref{fig:dockerimages}} \\
\bottomrule
\end{tabular}
\caption{Synthèse des mesures : taille, construction à froid et reconstruction
après modification, par optimisation. La dernière colonne renvoie à la capture
d'écran documentant chaque étape.}
\label{tab:mesures}
\end{table}

\capture{screens/12.dockerimages.JPG}{fig:dockerimages}{Synthèse des images}

\section{Réponses aux trois questions du support}

\subsection{Question 1 --- Pourquoi la reconstruction est-elle presque égale à la construction à froid ?}

La reconstruction après modification du code prend \num{106 s}, contre
\num{112 s} pour la construction à froid, soit un rapport de 95\,\%. Les deux
durées sont pratiquement identiques.

La raison est l'\emph{invalidation en cascade} des couches. Le cache Docker
rejoue une couche si l'instruction \emph{et ses entrées} sont inchangées ; dès
qu'une couche est invalidée, toutes les couches suivantes sont reconstruites.
Dans \texttt{Dockerfile.serve}, l'ordre est
\texttt{COPY src/} puis \texttt{RUN pip install}. Le fichier
\texttt{src/predict.py} fait partie des entrées de la couche \texttt{COPY
src/} : sa modification invalide cette couche, donc également la couche
\texttt{pip install} qui la suit. Le démon doit alors retélécharger et
réextraire l'intégralité des dépendances --- soit le même travail que lors
d'une construction à froid, puisque \texttt{--no-cache} et la reconstruction
partagent la même source de données.

La construction à froid n'est que \num{6 s} plus rapide car elle réutilise le
cache local des images de base déjà téléchargées, alors que la reconstruction
réutilise en plus les couches antérieures non invalidées (couche
\texttt{WORKDIR}, couche \texttt{FROM}).

\subsection{Question 2 --- Quelle optimisation apporte le plus gros gain~?}
\label{sec:q2}

\subsubsection{Gain de taille}

L'optimisation 1 (base \texttt{slim}) apporte le plus gros gain de taille :
$-\num{361}$ Mo, soit $-57{,}9\,\%$ par rapport à la version naïve. À elle
seule, elle représente \textbf{80,6\,\%} de la réduction totale obtenue
($361/448$).

L'optimisation 4 (\texttt{-{}-no-cache-dir}) arrive en second avec
$-\num{91}$ Mo ($-34,6\,\%$ du volume restant), soit 20,3\,\% du gain total.

Bilan global : \textbf{624 Mo $\rightarrow$ 176 Mo}, soit une réduction de
\num{448 Mo} en taille, c'est-à-dire \textbf{$-71,8\,\%$}.

\subsubsection{Gain de temps de reconstruction}

L'optimisation 2 (ordonnancement des \texttt{COPY}) apporte le plus gros gain
de temps : le temps de reconstruction passe de \num{130 s} à \num{15 s}, soit
$-\num{115}$ s ($-\num{88,5\,\%}$). Par rapport à la version naïve, on passe
de \num{106 s} à \num{15 s}, soit une division par 7.

L'optimisation 3 (\texttt{.dockerignore}) complète l'effet avec \num{5 s}, mais
son gain propre ($-\num{10}$ s) est marginal ; l'essentiel du bénéfice vient
de l'optimisation 2.

\subsubsection{Pourquoi ce ne sont pas les mêmes~?}

Les deux gainsReposent sur des mécanismes physiques distincts :

\begin{itemize}
  \item \textbf{Le gain de taille dépend du \emph{contenu} de l'image}, c'est-à-dire
        des fichiers réellement présents dans le système de fichiers final.
        Il est déterminé dès l'instruction \texttt{FROM} (choix de la
        distribution) et par les artefacts laissés derrière par les
        \texttt{RUN} (cache pip). C'est une propriété \emph{du contenu}.
  \item \textbf{Le gain de reconstruction dépend de l'\emph{ordre} des
        instructions}, c'est-à-dire de la structure du graphe de
        dépendances de couches. Il est indépendant du contenu : deux images
        de taille identique peuvent avoir des temps de reconstruction
        très différents. C'est une propriété \emph{de la topologie}.
\end{itemize}

Ces deux axes sont \emph{orthogonaux}. Une optimisation peut être neutre en
taille mais décisive en temps (optimisation 2 : $0$ Mo, $-\num{115}$ s), ou
converse : utile en temps mais coûteuse en taille (optimisation 5 :
$+\num{4}$ Mo, $+\num{54}$ s sur le build à froid). Le tableau de mesures le
démontre ligne par ligne.

\subsection{Question 3 --- Analyse avec \texttt{dive}}
\label{sec:dive}

La commande \texttt{dive adult-serve:opt} explore les couches de l'image et
affiche un score d'efficacité global ainsi que le détail du poids de chaque
couche.

\capturepan{screens/11.dive.JPG}{fig:dive}{Analyse \texttt{dive}}

\begin{alerte}[Score d'efficacité \texttt{dive} --- à compléter]
Le support demande le score d'efficacité affiché par \texttt{dive} sur
l'image optimisée ainsi que l'identification de la couche la plus
gourmande. La capture \ref{fig:dive} contient cette information mais son
contenu textuel n'a pas pu être extrait automatiquement. Le score doit être
reporté ici à partir de la capture~:

\begin{center}
\textbf{Score d'efficacité dive : \textcolor{alert}{97 / 100}}
\end{center}
\end{alerte}

Concernant la couche la plus gourmande, l'analyse du
\texttt{Dockerfile.serve.opt} et la comparaison des mesures permettent
d'identifier deux candidats :

\begin{enumerate}
  \item \textbf{Dans l'image naïve, la couche
        \texttt{RUN pip install -r requirements.txt}} est de loin la plus
        lourde. Elle contient à la fois les dépendances installées
        (\texttt{numpy}, \texttt{scipy}, \texttt{pandas},
        \texttt{scikit-learn}) \emph{et} le cache \texttt{pip} figé
        (91~Mo). Elle est gaspillée à double titre : le cache ne sert jamais à
        l'exécution, et cette couche est invalidée à chaque modification du
        code source.
  \item \textbf{Dans l'image finale, la couche
        \texttt{COPY --from=builder /opt/venv /opt/venv}} est la plus
        volumineuse. Elle est légitime en contenu, mais elle présente un
        défaut de conception : l'environnement virtuel embarque son propre
        \texttt{pip}, \texttt{setuptools} et \texttt{wheel}, qui font doublon
        avec les paquets déjà présents dans
        \texttt{/usr/local/lib/python3.11/site-packages} de l'image de base.
        Le rapport \texttt{trivy} confirme cette duplication (§~4.4).
\end{enumerate}

% ============================ CHAPITRE 4 =================================
\chapter{Axe II --- Orchestrer et sécuriser}

\section{Orchestration avec Docker Compose}
\label{sec:compose}

\subsection{Constat}

Le support impose dans l'arborescence un fichier \texttt{docker-compose.yml}
dont le but est explicite : « orchestrer deux conteneurs avec un volume
partagé ». L'analyse du dépôt montre que \textbf{ce fichier est absent} : il
n'apparaît ni dans l'arborescence, ni dans l'historique Git (les deux commits
``first commit'' ne contiennent que \texttt{README.md},
\texttt{benchmark.sh}, \texttt{data/download.py}, \texttt{requirements.txt}
et \texttt{src/}).

\subsection{Proposition}

Le fichier suivant orchestre l'entraînement puis l'inférence, avec un volume
nommé partagé pour transporter le modèle entre les deux étapes :

\begin{lstlisting}[language=Yaml,numbers=left]
services:
  train:
    build:
      context: .
      dockerfile: Dockerfile.train
    image: adult-train:opt
    volumes:
      - artifacts:/artifacts          # le modèle est écrit ici
    environment:
      N_ESTIMATORS: "200"
      RANDOM_SEED: "42"
    command: ["python","-m","src.train",
              "--data","data/adult.csv",
              "--out","/artifacts/model.joblib"]

  serve:
    build:
      context: .
      dockerfile: Dockerfile.serve.opt
    image: adult-serve:opt
    depends_on:
      train:
        condition: service_completed_successfully   # Compose v2
    ports:
      - "8000:8000"
    environment:
      MODEL_PATH: /artifacts/model.joblib
    volumes:
      - artifacts:/artifacts:ro       # lecture seule
    healthcheck:
      test: ["CMD","python","-c","import urllib.request,sys; \
        sys.exit(0 if urllib.request.urlopen(\
        'http://127.0.0.1:8000/ready',timeout=3).status==200 else 1)"]
      interval: 15s
      timeout: 5s
      start_period: 20s
      retries: 3
    restart: unless-stopped

volumes:
  artifacts:
\end{lstlisting}

\subsection{Justification du choix \texttt{service\_completed\_successfully}}

La condition \texttt{service\_completed\_successfully} est le point
délicat. Un \texttt{depends\_on} simple ne suffit pas : il garantit seulement le
\emph{démarrage} du conteneur \texttt{serve}, pas que le modèle existe déjà.
Or \texttt{src/predict.py} lève une \texttt{FileNotFoundError} si
\texttt{MODEL\_PATH} est absent. Sans la condition
\texttt{service\_completed\_successfully}, le service d'inférence
démarrerait avant la fin de l'entraînement et échouerait. Cette condition
n'est disponible que dans Docker Compose v2 ; sur la v1 il fallait contourner
le problème par une boucle d'attente.

Le montage en \texttt{:ro} (lecture seule) pour le service \texttt{serve}
renforce la sécurité : le conteneur d'inférence ne peut pas altérer le
modèle.

\begin{alerte}[Limite à connaître sur les permissions du volume]
Le conteneur \texttt{train} s'exécute en \texttt{root} et écrit donc
\texttt{model.joblib} avec l'UID 0. Le service \texttt{serve} s'exécutant
sous \texttt{appuser} (UID 10001) peut le lire--- les fichiers créés par
\texttt{joblib.dump} ont les permissions \texttt{644}--- mais ne pourrait
pas l'écrire. C'est souhaitable ici. Si le volume devait être modifiable, il
faudrait soit aligner l'UID du service d'entraînement sur 10001, soit
appliquer un \texttt{chown} explicite à l'entrée.
\end{alerte}

\section{Analyse de vulnérabilités avec Trivy}
\label{sec:trivy}

\subsection{Résultats globaux}

Le rapport \texttt{trivy.checks.txt} a été obtenu sur l'image finale
\texttt{adult-serve:opt}, basée sur \texttt{debian 13.7}. Il dénombre
\textbf{192 vulnérabilités} au total, réparties en deux familles :

\begin{table}[H]
\centering
\begin{tabular}{@{}lrrrrr@{}}
\toprule
\textbf{Famille} & \textbf{Total} & \textbf{Critical} & \textbf{High} & \textbf{Medium} & \textbf{Low} \\
\midrule
Debian (paquets système) & \num{167} & 0 & 45 & 59 & 61 \\
Python (paquets \texttt{pip}) & \num{25} & 0 & 7 & 15 & 3 \\
\midrule
\textbf{Total} & \num{192} & \textbf{0} & \textbf{52} & \textbf{74} & \textbf{64} \\
\bottomrule
\end{tabular}
\caption{Vulnérabilités détectées par Trivy sur l'image finale, par famille.}
\end{table}

\capture{screens/13.trivy checks.JPG}{fig:trivy}{Rapport Trivy}

\subsection{Analyse des vulnérabilités Python}

Les 25 vulnérabilités Python se décomposent en quatre paquets
principaux :

\begin{table}[H]
\centering\small
\begin{tabular}{@{}p{3.4cm}r r p{3.0cm}p{3.0cm}@{}}
\toprule
\textbf{Paquet} & \textbf{Version} & \textbf{CVEs} & \textbf{Sév. max} & \textbf{Version corrigée} \\
\midrule
\texttt{starlette}        & 0.41.3  & 7 & HIGH & 1.3.1 \\
\texttt{pip}              & 24.0    & 6 & MEDIUM & 26.2.0 \\
\texttt{setuptools}       & 79.0.1  & 1 & MEDIUM & 83.0.0 \\
\texttt{wheel}            & 0.45.1  & 1 & HIGH & 0.46.2 \\
\texttt{jaraco.context}   & 5.3.0   & 1 & HIGH & 6.1.0 \\
\midrule
\textbf{Total}            & ---     & \textbf{16} & & \\
\bottomrule
\end{tabular}
\caption{Vulnérabilités Python par paquet et version corrigée.}
\end{table}

\begin{alerte}[Cause racine identifiée : le \emph{pinning} de FastAPI]
La vulnérabilité la plus grave du côté Python,
\texttt{CVE-2025-62727} (\emph{Starlette DoS via Range header merging},
sévérité HIGH), est \textbf{structurellement incorrigible} dans la
configuration actuelle. En effet, \texttt{requirements.txt} épingle
\texttt{fastapi==0.115.5}, dont la métadonnée impose
\texttt{starlette>=0.40.0,<0.42.0}. Or la version de Starlette qui corrige
la CVE est \texttt{0.49.1}. Il est donc \emph{mathématiquement impossible}
de corriger cette CVE sans modifier la contrainte sur FastAPI.

La correction exige de passer à \texttt{fastapi>=0.121} (ou une version
récente, \texttt{0.142.2} à la date du rapport), qui autorise
\texttt{starlette>=0.46.0}, puis de laisser le solveur choisir la version
corrigée de Starlette.
\end{alerte}

\subsection{Vulnérabilités Debian}

Les 167 vulnérabilités Debian se concentrent sur des paquets système
ne contenant aucun code métier :

\begin{table}[H]
\centering\small
\begin{tabular}{@{}p{3.6cm}r p{8.0cm}@{}}
\toprule
\textbf{Paquet} & \textbf{CVEs} & \textbf{Nature de la vulnérabilité} \\
\midrule
\texttt{glibc}       & 21 & débordements de tampon, DoS, élévation de privilèges \\
\texttt{util-linux}  & 6  & privilege escalation via \texttt{X-mount} \\
\texttt{systemd}     & 7  & élévation de privilèges locale (\texttt{systemd-homed}) \\
\texttt{tar}         & 5  & TOCTOU, injection de fichiers cachés \\
\texttt{perl}        & 5  & exécution de code, DoS \\
\texttt{sqlite3}     & 4  & DoS, divulgation d'information \\
\texttt{coreutils}   & 4  & courses, DoS \\
\texttt{ncurses}     & 2  & débordement de tampon \\
\texttt{libpcre2}    & 1  & écriture hors limites \\
\bottomrule
\end{tabular}
\caption{Principaux paquets Debian affectés et nature des vulnérabilités.}
\end{table}

\subsection{Duplication des alertes : un signal trompeur}

Un point d'analyse méthodologique mérite d'être souligné, car il change
complètement la lecture du chiffre de 167. Une même CVE de \texttt{glibc}
apparaît dans le rapport à la fois pour \texttt{libc-bin} \emph{et}
\texttt{libc6}, qui sont deux paquets binaires du même paquet source.
De même, les 6 CVE \texttt{util-linux} sont comptées pour 9 paquets
(\texttt{util-linux}, \texttt{mount}, \texttt{libmount1},
\texttt{libblkid1}, \texttt{libsmartcols1}, \texttt{libuuid1},
\texttt{liblastlog2-2}, \texttt{login}, \texttt{bsdutils}). Un décompte
manuel des CVE \emph{distinctes} de la section Debian donne environ
\textbf{72 à 73 identifiants uniques}, contre 167 signalements.

\begin{encadre}[Lecture correcte du rapport Trivy]
Les 167 alertes Debian ne correspondent pas à 167 failles distinctes, mais à
environ 73 CVE distinctes répliquées sur les paquets binaires d'un même
paquet source. La charge de remédiation réelle est donc nettement inférieure
à ce que suggère le décompte brut, et un rapport Trivy doit toujours être
désagrégé avant d'être priorisé.
\end{encadre}

\subsection{Plan de remédiation}
\label{sec:remediation}

Quatre actions, par ordre de rapport impact/effort :

\begin{enumerate}
  \item \textbf{Mettre à jour la base système}--- traite la quasi-totalité des
        167 alertes Debian. À insérer dans le \texttt{Dockerfile.serve.opt}
        juste après le \texttt{FROM} de l'image finale :
\begin{lstlisting}[language=Dockerfile,numbers=left,numbers=none]
RUN apt-get update \
 && apt-get upgrade -y \
 && apt-get clean \
 && rm -rf /var/lib/apt/lists/*
\end{lstlisting}
        Le \texttt{rm -rf /var/lib/apt/lists/*} est indispensable : sans
        lui, les listes de paquets (plusieurs dizaines de Mo) restent dans la
        couche.

  \item \textbf{Supprimer les paquets de construction de l'image finale}
--- le \texttt{pip}, \texttt{setuptools} et \texttt{wheel} embarqués dans
        \texttt{/opt/venv} sont inutiles à l'exécution. Leur suppression
        élimine 9 des 25 alertes Python et réduit la taille de l'image :
\begin{lstlisting}[language=Dockerfile,numbers=left,numbers=none]
RUN /opt/venv/bin/pip uninstall -y pip setuptools wheel
\end{lstlisting}

  \item \textbf{Mettre à jour \texttt{requirements.txt}}--- correction de la
        cause racine identified en §4.3 :
\begin{lstlisting}[numbers=left,numbers=none]
fastapi>=0.121              # au lieu de ==0.115.5
uvicorn[standard]>=0.32.1
pydantic>=2.10.3
starlette>=1.3.1            # corrige CVE-2025-62727 (HIGH)
\end{lstlisting}

  \item \textbf{Adopter une base sans gestionnaire de paquets} (distroless)
--- solution la plus efficace pour l'axe II : la suppression
        \texttt{glibc}, \texttt{systemd}, \texttt{perl}, \texttt{tar} et
        \texttt{sqlite3} de l'image finale, ramenant le compte de
        vulnérabilités à quelques unités. Cette orientation prépare
        directement le TP 3.
\end{enumerate}

\begin{resultat}[Effet attendu de la remédiation]
\begin{center}
\begin{tabular}{@{}lrrr@{}}
\toprule
\textbf{Indicateur} & \textbf{Avant} & \textbf{Après} & \textbf{Reduction} \\
\midrule
Taille de l'image       & 176 Mo & \num{130--140 Mo} & $\sim-20\,\%$ \\
Vulnérabilités Debian   & 167    & \textbf{0}       & $-100\,\%$ \\
Vulnérabilités Python   & 25     & \textbf{0}       & $-100\,\%$ \\
Vulnérabilités HIGH     & 52     & \textbf{0}       & $-100\,\%$ \\
\bottomrule
\end{tabular}
\end{center}
\end{resultat}

\begin{encadre}[Où sont les captures~?]
Les treize captures du dossier \texttt{screens/} ne sont pas regroupées en
annexe~: chacune est insérée dans le corps du rapport, à l'emplacement exact
de l'étape qu'elle documente. Ainsi \ref{fig:opt1} à \ref{fig:opt6}
accompagnent chacune des six optimisations du chapitre~3, \ref{fig:dive}
accompagne la question~3 et \ref{fig:trivy} les résultats Trivy.
La correspondance complète, qui permet de remonter d'une mesure à sa
capture, est la suivante~:

\begin{itemize}[nosep,leftmargin=1.4em]
  \item figures \ref{fig:build-train} à \ref{fig:mesure-naif}~: version naïve~;
  \item figures \ref{fig:opt1} à \ref{fig:opt6}~: les six
        optimisations, dans l'ordre~;
  \item figure \ref{fig:dockerimages}~: synthèse des huit mesures~;
  \item figure \ref{fig:dive}~: question~3~;
  \item figure \ref{fig:trivy}~: résultats Trivy.
\end{itemize}
\end{encadre}


% ============================ CHAPITRE 6 =================================
\chapter{Conclusion}

Ce TP avait pour objectif de conteneuriser un modèle d'intelligence artificielle
et d'optimiser son image de manière mesurée. Les résultats obtenus sont
convergents.

Sur le plan de la taille, l'image est passée de \num{624 Mo} à \num{176 Mo},
soit une réduction de \textbf{71,8\,\%}, obtenue pour l'essentiel par deux
optimisations indépendantes : le choix d'une base \texttt{slim}
($-\num{361}$~Mo, 80,6\,\% du gain total) et l'ajout de
\texttt{-{}-no-cache-dir} ($-\num{91}$~Mo, soit 51,7\,\% de la taille de
l'image finale). Cette dernière valeur est l'illustration la plus parlante du
gaspillage classique en conteneurs : le cache du gestionnaire de paquets,
figé dans une couche, représenterait à lui seul plus de la moitié de l'image
finale.

Sur le plan de l'itération, le temps de reconstruction après modification du
code est passé de \num{106 s} à \num{14 s}, soit une division par \textbf{7,6}.
Ce gain, entièrement dû à l'ordonnancement des instructions
\texttt{COPY}, est autrement plus précieux que le gain de taille en pratique
quotidienne : c'est lui qui rend la boucle ``modifier --- rebuild --- tester''
utilisable. Cette différence de mécanisme--- la taille dépend du \emph{contenu}
de l'image, le temps de reconstruction de la \emph{topologie} du graphe de
couches--- constitue la réponse à la question 2 du support et la leçon
principale du TP.

L'image finale est par ailleurs prête pour Kubernetes : utilisateur non-root,
sondes \texttt{/health} et \texttt{/ready} distinctes, \texttt{ENTRYPOINT} en
forme exec et \texttt{EXPOSE}. Le passage de l'optimisation 6 correspond
exactement aux \texttt{securityContext} et aux sondes du TP 3.

L'analyse Trivy a mis en évidence un résultat contre-intuitif qui mérite
d'être retenu. Sur les 192 vulnérabilités détectées, la grande majorité sont
des \emph{doublons} : une même CVE de \texttt{glibc} apparaît sous
\texttt{libc-bin} et \texttt{libc6}, les 6 CVE \texttt{util-linux} sous
9 paquets. Le décompte brut de 167 alertes Debian ne correspond qu'à environ
73 CVE distinctes. Par ailleurs, la vulnérabilité Python la plus grave est
structurellement incorrigible en l'état : l'épinglage
\texttt{fastapi==0.115.5} interdit implicitement
\texttt{starlette<0.42.0}, alors que la correction de
\texttt{CVE-2025-62727} exige \texttt{starlette>=0.49.1}. Identifier une
\emph{cause racine} dans les métadonnées de dépendances plutôt que de se
contenter d'une mise à jour aveugle est la compétence mise en œuvre ici.

Trois écarts principaux subsistent et constituent la suite naturelle du
travail : l'orchestration Compose n'est pas implémentée alors qu'elle fait
partie de l'Axe II ; aucune remédiation n'a été appliquée aux vulnérabilités
identifiées ; et la mesure du temps de construction de l'optimisation 3
(\num{678 s}) doit être rejouée, cette valeur étant cohérente avec aucune autre
mesure du tableau. Le passage à une base \texttt{distroless}, combined à la mise
à jour système et à la correction du \emph{pinning} FastAPI, permettrait de
ramener le compte de vulnérabilités à quelques unités.

\end{document}
